% Generated draft: requires mathematical and source review.
\documentclass[11pt]{article}
\usepackage[margin=25mm]{geometry}
\usepackage{fontspec}
\setmainfont{Latin Modern Roman}
\usepackage{amsmath,amssymb,mathtools,unicode-math}
\setmathfont{Latin Modern Math}
\usepackage{enumitem,needspace,xurl,hyperref}
\setlength{\parindent}{0pt}
\setlength{\parskip}{5pt}
\setlength{\emergencystretch}{3em}
\newcommand{\N}{\mathbb{N}}
\newcommand{\Z}{\mathbb{Z}}
\newcommand{\Q}{\mathbb{Q}}
\newcommand{\R}{\mathbb{R}}
\newcommand{\C}{\mathbb{C}}
\newcommand{\F}{\mathbb{F}}
\newcommand{\A}{\mathbb{A}}
\newcommand{\PP}{\mathbb{P}}
\newcommand{\E}{\mathbb{E}}
\newcommand{\eps}{\varepsilon}
\newcommand{\dd}{\,\mathrm{d}}
\begin{document}
\section*{UnsolvedMath ID: 20003099}
\textbf{Source problem code:} catalogue entry
\subsection{Definitions and mathematical statement}
The source asks whether a Thurston geodesic has dense projection to moduli space, and recalls that Masur used closed Teichm\"uller geodesics to approximate a Teichm\"uller geodesic dense in moduli space. The supplied partial result is a dynamical criterion; it neither constructs such a Thurston geodesic nor proves the needed dynamical hypotheses.

Let \(P\) be a nonempty Baire phase space, \(M\) a second-countable base space, \(\pi:P\to M\) continuous, and \(\Phi_t:P\to P\), \(t\geq0\), a continuous semiflow. For a countable basis \((V_j)\) of nonempty open subsets of \(M\), define
\[
 U_j=\bigcup_{t\geq0}\Phi_t^{-1}(\pi^{-1}(V_j)),\qquad
 D=\{z\in P:\overline{\{\pi(\Phi_tz):t\geq0\}}=M\}
   =\bigcap_jU_j.
\]
The aid's base-visitation condition is
\[
 \forall\,O\subset P\text{ nonempty open},\quad
 \forall\,V\subset M\text{ nonempty open},\quad
 \exists z\in O\ \exists t\geq0:\ \pi(\Phi_tz)\in V.
\]
It asserts the equivalence between this condition and residuality of \(D\). Here residuality uses the Baire setting, so the resulting countable intersection of open dense sets is dense and nonempty. The quantifiers allow the intermediate point \(z\) and time to depend on \(O,V\); the Baire argument then yields points whose projected orbit meets every basis open set. This is base visitation, not full phase-space transitivity.

As an alternative sufficient input, the aid gives a full-support invariant ergodic probability measure on \(P\), with the relevant nonempty base opens having nonempty open preimages, as for a surjective projection. In that case positive measure of these preimages and ergodicity supply visitation, and hence dense projected forward orbits. The existence of this measure is an additional hypothesis, not established by the criterion.

For the reported geometric application, take a closed oriented surface \(S\) of genus at least two, its Teichm\"uller space \(\mathcal T(S)\), and moduli space \(\mathcal M(S)\). The aid refers to the phase quotient for either Pan--Wolf Thurston geodesic flow, with projection to \(\mathcal M(S)\). Under the stated continuity, Baire, and second-countability conditions, base visitation would produce a residual family of directed Thurston rays whose moduli projections are dense. The precise phase spaces and flow constructions are not supplied.

For the harmonic-stretch flow, the aid further states that the rays extend to directed Thurston lines. In a unit-speed lift \(\gamma:\R\to\mathcal T(S)\), directed means
\[
 d_{\mathrm{Th}}(\gamma(s),\gamma(t))=t-s\qquad(s<t),
\]
with the asymmetric order retained. A dense projected forward ray on such a line would answer the stated dense-projection question. No reverse-distance equality or globally minimizing geodesic in the quotient metric is asserted. Neither base visitation nor a full-support invariant ergodic probability measure for these flows is proved in the supplied aid, so this remains a conditional reduction of the original existence problem.
\subsection{Short English statement}
The supplied criterion reduces dense projected Thurston geodesics to a precise base-visitation property, or alternatively to an invariant ergodic measure with full support. A Baire argument would then give a residual set of dense projected rays, and harmonic-stretch extension would give lines. Neither required dynamical input is established.
\subsection{Sources}
\begin{itemize}
\item Local UnsolvedMath registry, permanent ID \(20003099\), exact source code \texttt{catalogue entry}; canonical \texttt{aim-topology-notes.json}, \texttt{notes[186]}, source tag \texttt{problem}.
\item AIM list ``Problems on Thurston Metric,'' workshop ``Lipschitz metric on Teichm\"uller space,'' Problem 2.9: \url{https://aimath.org/pastworkshops/lipschitzteichproblems.pdf}. The supplied verification identifies the canonical question as exact. This is a supplied citation only; no PDF was fetched.
\item The original text mentions Masur's approximation argument. The aid names Pan--Wolf geodesic flows and the harmonic-stretch flow, but supplies no paper titles, bibliographic identifiers, exact phase-space constructions, or full authoritative references. The original literature field is empty.
\item Research attempt 1: result \texttt{reduction}, problem status at run \texttt{open}, rigor label \texttt{proved}, novelty confidence \texttt{low}; classification \texttt{PARTIAL-PROGRESS}, registry status \texttt{partially\_solved}, suggested difficulty L4. Created and updated 14 August 2026. These are supplied assessments, not a proof of dense Thurston geodesic existence.
\item Attribution: Ulam AI / UnsolvedMath Contributors, CC BY 4.0. The machine-generated research aid is not peer reviewed and requires independent verification before citation.
\end{itemize}
\end{document}
